\section{Objetivos}

El objetivo principal de este trabajo es extender Voctomix hacia una plataforma de
código abierto, modular y reproducible para la producción remota de vídeo en tiempo real.
Partiendo de Voctomix como base, se busca construir un sistema completo que permita
realizar producciones audiovisuales de calidad profesional sin depender de \textit{hardware}
propietario de alto coste, y que pueda desplegarse fácilmente mediante contenedores.

Para lograrlo, se plantean los siguientes objetivos específicos:

\begin{enumerate}
  \item Ampliar las funcionalidades del mezclador con características propias de la
        producción profesional: modos de composición de imagen con transiciones animadas,
        rotulación dinámica con gráficos y textos en tiempo real,
        un gestor de estados de emisión y sincronización automática del audio con el vídeo.

  \item Implementar un servicio de telemetría que exporte el estado del sistema en
        formato \ac{JSON} y lo publique en el gestor de colas RabbitMQ, de forma que
        se pueda integrar con herramientas externas sin tocar el núcleo de la aplicación.

  \item Contenerizar el sistema completo con Docker, de modo que
        cualquiera pueda desplegarlo con un único comando y sin preocuparse por
        dependencias ni configuraciones del entorno.

  \item Extender el despliegue a Kubernetes, validando que el sistema funciona también
        en infraestructuras más complejas y multinodo.

  \item Validar el sistema en distintos escenarios de despliegue, desde un único
        ordenador hasta un clúster de Kubernetes, midiendo la latencia extremo a
        extremo, la estabilidad del \textit{pipeline} y el correcto funcionamiento del sistema
        de telemetría en cada configuración.

  \item Demostrar la aplicabilidad del sistema en un contexto de producción real,
        mediante su integración en el proyecto europeo CyberNEMO de la Universidad
        Politécnica de Madrid, donde opera como motor de producción audiovisual en
        un entorno de investigación activo.
\end{enumerate}
